\subsection{複数属性意思決定}

これまでの多くの技法は、金銭という単一基準で判断する。しかし、実務の判断では、金銭だけでは不十分である。セキュリティ、利用者満足、環境負荷、保守性、社会的責任、組織文化への適合性など、多くの属性が同時に関係する。

\subsubsection{補償型技法}

補償型技法は、複数の基準を一つの評価値へまとめる方法である。ある基準の低い点数を、別の基準の高い点数で補うことができる。代表例には、無次元化、加重和、階層分析法がある。

\begin{notebox}{例}
案 A は費用が高いが信頼性が非常に高い。案 B は費用が低いが信頼性が低い。補償型技法では、費用と信頼性に重みを付け、総合点で比較できる。
\end{notebox}

\subsubsection{非補償型技法}

非補償型技法は、基準間の代替を認めない方法である。たとえば「安全基準を満たさない案は、どれほど安くても採用しない」といった判断である。代表例には、支配、満足化、辞書式順序がある。

\begin{notebox}{用語の見分け方}
重大な安全性、法令順守、倫理的制約は、補償型で扱うと危険である。費用が安いから安全基準未達を許す、という判断は成立しない場合が多い。
\end{notebox}
